% GENERATED by data/scripts/eval/gen_component_cost.py -- do not edit.
% Per-circuit cost profile for Q2 (component R1CS per folding step).
% Each c_i cell is `raw (adjusted)': raw = component R1CS / input byte;
% adjusted = raw after the MEASURED c4 log-up query cost (PERF 1012)
% is moved out of the framework and split across tiers by data-column
% proportion, so c1/c2/c3 become all-in and c4 keeps only Poseidon/IO.
% A `*' row lacks the per-component data-column log, so only raw is
% shown (the split is not yet available -- not accurate).
\begin{table*}[t]
\centering
\small
\setlength{\tabcolsep}{6pt}
\begin{tabular}{@{}ll r r rrrr r@{}}
\toprule
        &      & input & total  & \multicolumn{4}{c}{R1CS / input byte: raw (adjusted)}
        & corpus \\
\cmidrule(lr){5-8}
Dataset & Circ & (B)   & R1CS   & $c_1$ (CP) & $c_2$ (SDE) & $c_3$ (DFA)
        & $c_4$ (frwk) & share \\
\midrule
\textsc{Mal} & *** & *** & *** & *** & *** & *** & *** & *** \\
\midrule
\textsc{Dna} & *** & *** & *** & *** & *** & *** & *** & *** \\
\midrule
\textsc{Dlp} & *** & *** & *** & *** & *** & *** & *** & *** \\
\bottomrule
\end{tabular}
\caption{Per-circuit cost profile. Each dataset is discharged by a small set of circuits. \emph{input} is the bytes processed per folding step. Each $c_i \in (c_1,c_2,c_3,c_4)$ is the per-input-byte R1CS constraints of the four components of a circuit.Each $c_i$ is shown as \emph{raw\,(adjusted)}: the \emph{adjusted} value moves the \emph{measured} log-up query cost in $c_4$ and charges it to each tier in  proportion to its data-column size, so $c_1/c_2/c_3$ reflect the  all-in per-byte cost and $c_4$ retains only folding overhead  and the processing of Logup share of lookup table. }\label{tbl:component-cost}
\end{table*}
